\documentclass[a4paper]{article}

% Layout
\usepackage[margin=1in]{geometry}

\usepackage{microtype}
\usepackage{enumitem}
\usepackage[toc,page]{appendix}

% Mathematics
\usepackage{amsmath,amssymb,amsfonts,amsthm,mathtools}
\usepackage{mathrsfs}
\usepackage{physics}
\usepackage{cancel}

% Graphics
\usepackage{xcolor}
\usepackage{tikz}
\usepackage{tikz-cd}
\usepackage{pgfplots}
\pgfplotsset{compat=1.18}

% Fonts
\usepackage[T1]{fontenc}
\usepackage[mathscr]{euscript}
\usepackage{mathbbol}
\usepackage[lite]{mtpro2}

% Paragraph typography
\setlength{\parskip}{0pt}
\setlength{\parindent}{1.25em}

% Utilities
\usepackage{etoolbox}
\usepackage{dynkin-diagrams}

% Warnings
\usepackage[l2tabu,orthodox]{nag}

% Environments
\newtheorem{theorem}{Theorem}[section]
\newtheorem{lemma}[theorem]{Lemma}
\newtheorem{proposition}[theorem]{Proposition}
\newtheorem{corollary}[theorem]{Corollary}
\newtheorem{conjecture}{Conjecture}

\theoremstyle{definition}
\newtheorem{definition}[theorem]{Definition}
\newtheorem{example}[theorem]{Example}

\theoremstyle{remark}
\newtheorem{remark}[theorem]{Remark}
\newtheorem{outline}[theorem]{Outline}

% Dynkin Diagrams
\def\row#1/#2!{\dynkin{#1}{#2}\\}

\newcommand{\tble}[1]{
   \renewcommand*\do[1]{\row##1!}
   \[
      \begin{array}{ll}
         \docsvlist{#1}
      \end{array}
   \]
}

% Extra Fonts
\DeclareMathAlphabet{\mathpzc}{OT1}{pzc}{m}{it}

\newcommand{\goth}[1]{\mathfrak{#1}}
\newcommand{\pzc}[1]{\mathpzc{#1}}

% Commands for FONT characters
\let\SS\S
\let\ag\gg

\ExplSyntaxOn

% Goth Font
\clist_map_inline:nn
{a,b,c,d,e,f,g,h,i,j,k,l,m,n,o,p,q,r,s,t,u,v,w,x,y,z,A,B,C,D,E,F,G,H,I,J,K,L,M,N,O,P,Q,R,S,T,U,V,W,X,Y,Z}
{
   \cs_undefine:c { g#1 }
   \cs_new:cpn { g#1 } { \goth{#1} }
}

% Scr Font
\clist_map_inline:nn
{A,B,C,D,E,F,G,H,I,J,K,L,M,N,O,P,Q,R,S,T,U,V,W,X,Y,Z}
{
   \cs_new:cpn { f#1 } { \mathscr{#1} }
}

% Cal Font
\clist_map_inline:nn
{A,B,C,D,E,F,G,H,I,J,K,L,M,N,O,P,Q,R,S,T,U,V,W,X,Y,Z}
{
   \cs_new:cpn { s#1 } { \mathcal{#1} }
}

% mathbf Font
\clist_map_inline:nn
{A,B,C,D,E,F,G,H,I,J,K,L,M,N,O,P,Q,R,S,T,U,V,W,X,Y,Z}
{
   \cs_undefine:c { #1 }
   \cs_new:cpn { q#1 } { \mathbf{#1} }
}

% BB Font
\clist_map_inline:nn
{A,B,C,D,E,F,G,H,I,J,K,L,M,N,O,P,Q,R,S,T,U,V,W,X,Y,Z}
{
   \cs_undefine:c { #1 }
   \cs_new:cpn { #1 } { \mathbb{#1} }
}

\ExplSyntaxOff

% Categories
\newcommand{\cC}{\mathbf{C}}
\newcommand{\cSet}{\mathbf{Set}}
\newcommand{\cAb}{\mathbf{Ab}}
\newcommand{\cGroup}{\mathbf{Grp}}
\newcommand{\ckSch}[1]{\mathbf{Sch}(#1)}
\newcommand{\ckAlg}[1]{\mathbf{Alg}(#1)}
\newcommand{\cMod}[1]{\mathbf{Mod}(#1)}
\newcommand{\cdh}{h}
\newcommand{\opp}[1]{{#1}^\text{op}}

\newcommand{\plat}[1]{{#1}_{\text{fl}}}
\newcommand{\etale}[1]{{#1}_\text{ét}}
\newcommand{\zar}[1]{{#1}_\text{Zar}}

\newcommand{\et}[1]{{#1}_\text{ét}}
\newcommand{\sep}[1]{{#1}^\text{sep}}
\newcommand{\mgbar}[1]{\overline{\fM}_{#1}}

\DeclareMathOperator{\cSch}{\mathbf{Sch}}

\DeclareMathOperator{\Mor}{Mor}

% Relative stuff
\DeclareMathOperator{\rSpec}{\mathpzc{Spec}}
\DeclareMathOperator{\rHom}{\mathpzc{Hom}}
\DeclareMathOperator{\rEnd}{\mathpzc{End}}
\DeclareMathOperator{\rProj}{\mathpzc{Proj}}
\DeclareMathOperator{\rExt}{\mathpzc{Ext}}

\DeclareMathOperator{\edim}{edim}
\DeclareMathOperator{\res}{res}

% Linear Algebra/Abelian Category/General Notions (?)
\DeclareMathOperator{\rk}{rank}
\DeclareMathOperator{\im}{im}
\DeclareMathOperator{\coker}{coker}
\DeclareMathOperator{\End}{End}
\DeclareMathOperator{\Hom}{Hom}
\DeclareMathOperator{\Aut}{Aut}
\DeclareMathOperator{\len}{length}
\DeclareMathOperator{\height}{height}
\DeclareMathOperator{\codim}{codim}
\DeclareMathOperator{\id}{id}

% Intersection
\DeclareMathOperator{\chern}{ch}
\DeclareMathOperator{\todd}{td}

% AG/Commutative Algebra/Algebra
\DeclareMathOperator{\Hilb}{Hilb}
\DeclareMathOperator{\Def}{Def}
\DeclareMathOperator{\car}{char}
\DeclareMathOperator{\proj}{Proj}
\DeclareMathOperator{\spec}{Spec}
\DeclareMathOperator{\pic}{Pic}
\DeclareMathOperator{\cacl}{CaCl}
\DeclareMathOperator{\cl}{Cl}
\DeclareMathOperator{\supp}{Supp}
\DeclareMathOperator{\Gr}{Gr}
\DeclareMathOperator{\Div}{Div}
\DeclareMathOperator{\Der}{Der}
\DeclareMathOperator{\Num}{Num}
\DeclareMathOperator{\Ann}{Ann}
\DeclareMathOperator{\Syl}{Syl}
\DeclareMathOperator{\Sym}{Sym}
\DeclareMathOperator{\Stab}{Stab}
\DeclareMathOperator{\Perm}{Perm}
\DeclareMathOperator{\Orb}{Orb}
\DeclareMathOperator{\Gal}{Gal}
\DeclareMathOperator{\Supp}{Supp}
\DeclareMathOperator{\trd}{tr.d.}
\DeclareMathOperator{\Frac}{Frac}
\DeclareMathOperator{\elm}{elm}

\newcommand{\red}[1]{{#1}_\text{red}}
\newcommand{\gsecx}[1]{\Gamma(X, #1)}
\newcommand{\sh}[1]{{#1}^\text{sh}}

% Homological Algebra
\DeclareMathOperator{\Ext}{Ext}
\DeclareMathOperator{\Tor}{Tor}
\DeclareMathOperator{\hd}{hd}
\DeclareMathOperator{\pd}{pd}
\DeclareMathOperator{\cH}{H}
\DeclareMathOperator{\Ti}{T}
\DeclareMathOperator{\dR}{R}
\DeclareMathOperator{\depth}{depth}

% Linear Algebraic Groups
\DeclareMathOperator{\PGL}{PGL}
\DeclareMathOperator{\SO}{SO}
\DeclareMathOperator{\SU}{SU}
\DeclareMathOperator{\Sp}{Sp}
\DeclareMathOperator{\GL}{GL}
\DeclareMathOperator{\Orth}{O}
\DeclareMathOperator{\Uni}{U}

% Random Stuff
\DeclareMathOperator{\Hasse}{Hasse}
\newcommand{\dlog}{\dd \log}

% Miscellaneous
\author{James Lee}

\allowdisplaybreaks

% Solarized
% \pagecolor[RGB]{0,20,26}
% \color[RGB]{191,191,191}

% Sephia
% \pagecolor[RGB]{249,239,220}

% Grey
% \pagecolor[RGB]{200, 200, 200}

% Black
% \pagecolor[RGB]{24,24,24}
% \color[RGB]{170,180,190}

% Cyber Deep Sea
% \pagecolor[RGB]{10,18,24}
% \color[RGB]{176,204,208}

\newcommand\restr[2]{{\left.\kern-\nulldelimiterspace #1 \vphantom{\big|} \right|_{#2}}}

\let\oldbigoplus\bigoplus

\renewcommand{\bigoplus}{\mathop{\textstyle\oldbigoplus}\displaylimits}


\title{Chapter 5, Section 5}

\begin{document}

\maketitle

\begin{enumerate} [label=\textbf{\arabic*.}, leftmargin=0em]

    \item Let $f$ be a rational function on the surface $X$.
          Show that it is possible to ``resolve the singularities of $f$'' in the following sense: there is a birational morphism $g : X' \to X$ so that $f$ induces a morphism of $X'$ to $\P^1$.

          \begin{proof}
              Assume $f$ is nonconstant, and let $f : X \dashrightarrow \P^1$ be the rational map induced by $f$.
              By (I, Ex. 4.2), there is a largest open set $U \subset X$ on which $f$ is represented by a morphism $\varphi : U \to \P^1$.
              If $t$ is a point in $\P^1$, denote by $f^{-1}(t)$ the closure of $\varphi^{-1}(t)$ in $X$.
              The proof of (5.1) follows through for any rational map of smooth projective varieties.
              In particular, the locus of indeterminacy $Z = X - U$ of $f$ is a finite set of points, and the fibers $\varphi^{-1}(t) \subset U$ are nonempty of dimension 1 (II, Ex. 3.22).
              Indeed, if $\varphi^{-1}(t)$ is empty for some $t$, then $\varphi$ factors through $\A^1 \hookrightarrow \P^1$, which implies the existence of a regular function on $U$.
              This is not possible, because $\dim(X - U) = 0$, so every regular function on $U$ extends uniquely to a regular function on $X$, and since $X$ is projective, the only regular functions on $X$ are the constant ones.

              Now, we claim that if $P$ is a closed point in $X$, then $P$ is a point of indeterminacy of $f$ if and only if $P$ is contained in two distinct fibers of $f$.
              The reverse direction is clear by continuity, and conversely, if $P$ is a point contained in only one fiber $f^{-1}(t)$, then the map $\varphi' : U \cup \{ P \} \to \P^1$ defined by $\restr{\varphi'}{U} = \varphi$ and $\varphi'(P) = t$ is also a morphism representing $f$, so by assumption on $U$ being the largest open set for which $f$ is a morphism, $P$ must be in $U$.
              Write the divisor of $f$ as $(f) = \sum n_i C_i$.
              Applying the embedded resolution (3.9) to the curve $Y = \bigcup C_i$ gives a birational morphism $\pi : \tilde{X} \to X$ such that $\pi^{-1}(Y)$ is a divisor with normal crossings.
              Let $\tilde{f} = f \circ \pi$, and let $\tilde{Z} \subset \tilde{X}$ be the locus of indeterminacy of $\tilde{f}$, which is also a finite set of points contained in $\pi^{-1}(Z)$.
              Similarly define the fiber $\tilde{f}^{-1}(t) = \pi^{-1}(f^{-1}(t))$.
              The divisor of $\tilde{f}$ can be written as a difference of effective divisors $(\tilde{f}) = E - D$, where $E = \sum a_i E_i$, $D = \sum b_j D_j$, every summand is an irreducible nonsingular curve, and whenever $E_i, D_j$ intersect at a point $P \in \tilde{X}$, the local equations form a part of a regular system of parameters at $P$.
              The latter provision combined with (Ex. 3.2) implies that the strict transforms of $E_i$ and $D_j$ are disjoint after a sufficient number of blow ups of $\tilde{X}$ at $P$, and for the same reason as above, a point of intersection of $E_i, D_j$ must be a point in $\tilde{Z}$.
              Let $\pi' : X' \to \tilde{X}$ be a composition of monoidal transformations so that $\pi'^{-1}(E)$ and $\pi'^{-1}(D)$ are disjoint, and let $f' = \tilde{f} \circ \pi'$.
              If $\tilde{E}$ and $\tilde{D}$ are disjoint, then $\tilde{E}.\tilde{D} = 0$, and since any two fibers of $\tilde{f}$ are linearly equivalent with a possible intersection at a point in $\tilde{Z}$, the strict transforms of every fiber $\tilde{f}^{-1}(t)$ in $X'$ are also disjoint.
              Hence, $f'$ is a morphism of $X'$ to $\P^1$.
          \end{proof}

          % \item Let $Y \cong \P^1$ be a curve in a surface $X$, with $Y^2 < 0$.
          % Show that $Y$ is contractible (5.7.2) to a point on a projective variety $X_0$ (in general singular).

          % \item If $\pi : \tilde{X} \to X$ is a monoidal transformation with center $P$, show that $\cH^1(\tilde{X}, \Omega_{\tilde{X}}) \cong \cH^1(X, \Omega_X) \oplus k$.
          % This gives another proof of (5.8).

          % \item Let $f : X \to X'$ be a birational morphism of nonsingular surfaces.
          % \begin{itemize}
          %     \item[(a)] If $Y \subseteq X$ is an irreducible curve such that $f(Y)$ is a point, then $Y \cong \P^1$ and $Y^2 < 0$.
          %     \item[(b)] (Mumford [6].) Let $P' \in X'$ be a fundamental point of $f^{-1}$, and let $Y_1, \dots, Y_r$ be the irreducible components of $f^{-1}(P')$.
          %     Show that the matrix $||Y_i.Y_j||$ is negative definite.
          % \end{itemize}

    \item[\textbf{5.}] Let $C$ be a curve, and let $\pi : X \to C$ and $\pi' : X' \to C$ be two geometrically ruled surfaces over $C$.
          Show that there is a finite sequence of elementary transformations (5.7.1) which transform $X$ into $X'$.

          \begin{proof}
              Let $P$ be a point in $X$, and let $L$ be the fiber of $\pi$ containing $P$.
              Consider the elementary transform of $X$ with center at $P$, which we denote by $\elm_P X$.
              Consider the diagram
              \[ \begin{tikzcd}
                      & \tilde{X} \arrow[ld, "\pi_1"'] \arrow[rd, "\pi_2"] &          \\
                      X &                                                    & \elm_P X
                  \end{tikzcd}, \]
              where $\pi_1$ is the monoidal transformation of $X$ centered at $P$ with exceptional curve $E$, and if $\tilde{L}$ is the strict transform of $L$ on $\tilde{X}$, then $\pi_2$ is the contraction of $\tilde{L}$, i.e., $Q = \pi_2(\tilde{L})$ is a point.
              Let $M = \pi_2(E)$, and let $D$ be a section of $\pi$ containing $P$.
              By (3.6), if $\tilde{D}$ is the strict transform of $D$ on $\tilde{X}$, then $\tilde{D} \sim \pi_1^*D - E$, so if $D'$ is the strict transform of $D$ by $\elm_P$, then $D' \sim \pi_{2*}\tilde{D}$.
              Also, $\pi_2$ is the monoidal transformation of $\elm_P X$ centered at $Q$, and $D'$ does not meet $Q$, so $\pi_2^*D' = \tilde{D}$ by (3.6).
              Thus by (3.2), $D'^2 = \tilde{D}^2 = D^2 - 1.$

              If $X$ has invariant $e$, then it is defined by the projective bundle $\P(\fE)$, where $\fE$ is a normalized locally free sheaf of rank 2, and $\bigwedge^2 \fE \cong \fO_C(\ge)$ for some divisor $\ge$ of degree $-e$ (2.9).
              Let $C_0$ be the section of $\pi$ corresponding to the surjection $\fE \to \fO_C(\ge)$.
              If $D \sim C_0 + (\gd - \ge)f$ is any section of $\pi$, then $\deg(\gd - \ge) \geq 0$ by the assumption that $\fE$ is normalized.
              A consequence is $e \geq - D^2$, or expressed differently, $-e$ is the minimum value for the self-intersection of a section on a ruled surface.
              Let $X_1$ be an elementary transform of $X$ centered at a point in $C_0$, and let $\tilde{C}_0$ be the strict transform of $C_0$.
              If $X_1$ has invariant $e_1$, then it follows by the above that $e_1 \geq -\tilde{C}_0^2 = e + 1$.
              Therefore, $X$ can be transformed into a geometrically ruled surface with arbitrarily high invariant.
              By (2.12), $\fE$ is decomposable if $e > 2g - 2$, so it remains to prove the statement for geometrically ruled surfaces determined by decomposable locally free sheaves.
              In fact, it suffices to show $X$ can be transformed to the trivial ruled surface $\P(\fO_C \oplus \fO_C) \cong C \times \P^1$.

              Now, let us consider how $\fE$ behaves under elementary transformations of $X$ if $\fE = \fO_C \oplus \fO_C(\ge)$ and $e = -\deg{\ge} > 0$.
              Ruled surfaces determined by decomposable locally free sheaves have two distinguished sections $C_0, C_1$ corresponding to the surjections of $\fE$ onto each summand $\fO_C(\ge), \fO_C$, respectively, such that $C_0^2 = -e$, $C_1^2 = e$, and $C_0 \cap C_1 = \emptyset$.
              Let $P$ be a point in $X$ lying over $R = \pi(P) \in C$, and consider the elementary transform $X''$ of $X$ centered at $P$.
              Following the notation of the above, we have a commutative diagram
              \[ \begin{tikzcd}
                      & \tilde{X} \arrow[ld, "\pi_1"']                     &                                       \\
                      X &                                                    & X'' \arrow[lu, "\pi_2^{-1}"', dashed] \\
                      & C \arrow[ru, "\sigma_i''"'] \arrow[lu, "\sigma_i"] &
                  \end{tikzcd} , \]
              where $\sigma_i, \sigma_i''$ are sections defining either $C_i, \tilde{C}_i$ for each $i = 0, 1$.
              Let $\tilde{C}_0, \tilde{C}_1$ denote the strict transforms of $C_0, C_1$, respectively.
              There are two cases we will consider, which are $P \in C_0$ or $P \in C_1$, and in either case, $\tilde{C}_0 \cap \tilde{C}_1 = \emptyset$, so $X''$ is decomposable by (Ex. 2.2).

              \fbox{$P \in C_0$} We claim $X'' \cong \P(\fO_C \oplus\fO_C(\ge - R))$.
              Under the identifications $\tilde{C}_i \cong C$ for each $i = 0, 1$ via $\sigma_i''$, we have the isomorphisms
              \begin{equation*}
                  \fO_{X''}(\tilde{C}_0) \otimes \fO_{\tilde{C}_0} \cong \fO_C(\ge - R), \quad \fO_{X''}(\tilde{C}_1) \otimes \fO_{\tilde{C}_1} \cong \fO_C.
              \end{equation*}
              Indeed, $\pi_2^{-1} \circ \sigma_0''$ is a morphism, and $\pi_2^{-1}$ is an isomorphism outside of the strict transform of the fiber containing $P$, so there is an isomorphism of invertible sheaves under the identification $\tilde{C}_0 \cong C$,
              \begin{equation*}
                  \fO_{X''}(\tilde{C}_0)  \otimes \fO_{\tilde{C}_0} \cong \sigma_0''^* \fO_{X''}(\tilde{C}_0) \cong (\pi_2^{-1} \circ \sigma_0'')^* \fO_{\tilde{X}}(\pi_1^*C_0 - E) \cong \fO_C(\ge - R),
              \end{equation*}
              and similarly for $\fO_{X''}(\tilde{C}_1) \otimes \fO_{\tilde{C}_1} \cong \fO_C$.
              We conclude by the proof of (Ex. 2.2) that $X'' \cong \P(\fE'')$, where $\fE'' = \fO_C \oplus \fO_C(\ge - R)$.

              \fbox{$P \in C_1$} We claim $X'' \cong \P(\fO_C \oplus \fO_C(\ge + R))$.
              Similar to the previous case, we have the isomorphisms
              \begin{equation*}
                  \fO_{X''}(\tilde{C}_0) \otimes \fO_{\tilde{C}_0} \cong \fO_C(\ge), \quad \fO_{X''}(\tilde{C}_1) \otimes \fO_{\tilde{C}_1}\cong \fO_C(-R).
              \end{equation*}
              Therefore, $X'' \cong \P(\fE'')$, where $\fE'' = \fO_C(-R) \oplus \fO_C(\ge)$, which is isomorphic to the projective bundle of $\fO_C \oplus \fO_C(\ge + R)$.

              If we write $\ge = \sum n_i P_i$, performing the appropriate elementary transforms centered at points in the fibers over $P_i$ will give a ruled surface $X'' = \P(\fE'')$ with $\fE'' = \fO_C \oplus \fO_C$.
              Hence, any ruled surface over $C$ can be transformed to $C \times \P^1$ with a finite sequence of elementary transforms.
          \end{proof}

          % \item[\textbf{6.}] Let $X$ be a surface with function field $K$.
          %       Show that every valuation ring $R$ of $K/k$ is one of the three kinds described in (II, Ex 4.12).

          %       \begin{proof}

          %       \end{proof}

          % \item Let $Y$ be an irreducible curve on a surface $X$, and suppose there is a morphism $f : X \to X_0$ to a projective variety $X_0$ of dimension 2, such that $f(Y)$ is a point $P$ and $f^{-1}(P) = Y$.
          % Then show that $Y^2 < 0$.

          % \item \textit{A surface singularity.} Let $k$ be an algebraically closed field, and let $X$ be the surface in $\A_k^3$ defined by the equation $x^2 + y^3 + z^5 = 0$.
          % It has an isolated singularity at the origin $P = (0, 0, 0)$.
          % \begin{itemize}
          %     \item[(a)] Show that the affine ring $A = k[x, y, z] / (x^2 + y^3 + z^5)$ of $X$ is a unique factorization domain, as follows.
          %     Let $t = z^{-1}$, $u = t^3x$, and $v = t^2 y$.
          %     Show that $z$ is irreducible in $A$, $t \in k[u, v]$, and $A[z^{-1}] = k[u, v, t^{-1}]$.
          %     Conclude that $A$ is a UFD.
          %     \item[(b)] Show that the singularity at $P$ can be resolved by eight successive blowing-up.
          %     If $\tilde{X}$ is the resulting nonsingular surface, then the inverse image of $P$ is a union of eight projective lines, which intersect each other according to the Dynkin diagram $\qE_8$:
          %     \tble{E/{8}}
          %     Here each circle denotes a line, and two circles are joined by a line segment whenever the corresponding lines intersect.
          % \end{itemize}

\end{enumerate}

\end{document}
